% ==========================================================================
% ABSTRACT
%
% EEE requirement: maximum 200 words, on a single page.
% Written as ONE flowing paragraph answering: problem/objective, why it
% matters, methods, key results (quantitative), and implications.
% ==========================================================================

\prelimchapter{Abstract}

Cardiovascular diseases are among the leading causes of death worldwide,
and timely detection of cardiac arrhythmias from the electrocardiogram
(ECG) is essential to prevent life-threatening complications. Manual ECG
interpretation is time-consuming, requires expert cardiologists, and is
vulnerable to inter-observer variability, which motivates automated
computer-aided classification. This thesis develops and compares two
deep-learning architectures for ECG beat classification: a hybrid
CNN--BiLSTM model with temporal additive attention, and a lightweight
multi-scale feature-augmented network (MSFT-Net) with squeeze-and-excitation
gating and FiLM conditioning. Both models classify the five ANSI/AAMI EC57
beat classes (N, S, V, F, Q) from one-second R-peak-centred windows combined
with rhythm, morphological, and template-similarity features under a
record-disjoint, inter-patient evaluation protocol, together with a
permissive intra-patient six-category ECG comparison that additionally
includes a CHF-derived clinical-condition category. Under the inter-patient
protocol, MSFT-Net achieved 89.45\% accuracy and 87.25\% macro-F1 with
167{,}090 parameters, outperforming the larger recurrent hybrid (82.00\%
accuracy, 73.26\% macro-F1), while the intra-patient protocol reached 98.59\%
accuracy. The INT8-quantised MSFT-Net was deployed on an STM32
microcontroller with 185.06~ms per-beat inference latency, demonstrating
that accurate, explainable, inter-patient arrhythmia classification is
feasible on resource-constrained wearable hardware.

\vspace{18pt}

\noindent
\textbf{Keywords:} ECG, arrhythmia classification, deep learning,
MSFT-Net, CNN--BiLSTM--attention, inter-patient evaluation, XAI, Grad-CAM,
microcontroller deployment

\clearpage